%\begin{figure} [h]
	\centering
	\pgfplotstableread[col sep=semicolon] {misc/chap_4_ball_bouncing/target_height_significance_post_hoc.csv}\datatable
	\tikzsetnextfilename{target_height_significance_post_hoc}
	\begin{tikzpicture}
		
		\begin{axis}[%
			height=0.3\linewidth,
			scale only axis,
			%xmin=0.65,
			%xmax=2.35,
			%xtick={1,1.5,2},
			xtick={1.5,1.75,2},
			xticklabels={{$h_1^*$},{$h_2^*$},{$h_3^*$}},
			%ymin=1.06276752948761,
			%ymax=1.78985784053802,
			ylabel=$\sigma(\epsilon_r)$ (\si{\centi\meter}),
			ylabel near ticks,
			ylabel style={yshift=-1mm},
			%boxplot/draw direction=y, 
			]
			\addplot [color=Orient] table[x=height, y=mean]{\datatable};
			\addplot [color=Orient, only marks, mark=*, mark size=1.7pt, mark options={fill=white, fill opacity=1}]
			plot [error bars/.cd, y dir = both, y explicit]
			table[x=height, y=mean, y error=SEM]{\datatable};		
		\end{axis}		
	\end{tikzpicture}%
	%\caption{Effets des variations de hauteur de la cible sur $\sigma(\epsilon_r)$, l'erreur type est utilisée comme intervalle de confiance.}
	\caption[ANOVA: Effets de la hauteur cible sur l'écart type des erreurs de rebond]{Effets de $h$ sur $\sigma(\epsilon_r)$.}
	\label{fig_target_height_significance_post_hoc}
%\end{figure} 